\subsection{TUBES AND TRANSVERSALS}

Lemma 4. Let H be a compact Lie group, \(\rho : \mathrm{H} \to \mathbf{GL}(\mathrm{V})\) a continuous (hence analytic) representation of H on a finite dimensional real vector space, and W a neighbourhood of the origin in V. There exists an open neighbourhood B of the origin, contained in W and stable under H, and an analytic isomorphism \(u: \mathrm{V} \to \mathrm{B}\) , compatible with the operations of H, such that \(u(0) = 0\) and \(\mathrm{Du}(0) = \mathrm{Id}_{\mathrm{V}}\) .

Choose a scalar product on V invariant under H (§1, no. 1). There exists a real number \(r > 0\) such that the open ball B of radius \(r\) is contained in W; it is clearly stable under H. Put \(u(v) = r(r^2 + \|v\|^2)^{-1/2} v\) for all \(v \in V\) ; then \(u\) is a bijective analytic map from V to B, compatible with the operations of H, and its inverse map \(w \mapsto r(r^2 - \|w\|^2)^{-1/2} w\) is analytic. Moreover, \(u(0) = 0\) and \(\mathrm{D}u(0) = \mathrm{Id}_{\mathrm{V}}\) .

PROPOSITION 5. Let H be a compact Lie group, \((h,x)\mapsto hx\) a law of left operation of class \(C^{r}\) of H on X, and x a point of X fixed under the operation of H. Then the group H operates linearly on the vector space \(T=T_{x}(X)\) ; there exists an open embedding \(\varphi:T\to X\) of class \(C^{r}\) , compatible with the operations of H, such that \(\varphi(0)=x\) and \(T_{0}(\varphi)\) is the identity map of T.

Let \((\mathrm{U},\psi,\mathrm{E})\) be a chart of X at x, such that U is stable under H (no. 2, Lemma 3) and such that \(\psi(x)=0\) . Identify E with T by means of \(\mathrm{T}_{x}(\psi)\) , and put

\[
\psi^ {\sharp} (y) = \int_ {\mathrm{H}} h. \psi (h ^ {- 1} y) d h \quad \text { for } y \in \mathrm{U},
\]

where dh is the Haar measure on H of total mass 1.

Then (§6, no. 4, Cor. 1) \(\psi^{\sharp}\) is a morphism of class \(\mathrm{C}^r\) from U to T, compatible with the operations of H, such that \(\psi^{\sharp}(x) = 0\) and \(d_x\psi^\sharp = \mathrm{Id}_\mathrm{T}\) . Hence, there exists an open set \(\mathrm{U}' \subset \mathrm{U}\) containing \(x\) , and an open neighbourhood V of 0 in T, such that \(\psi^{\sharp}\) induces an isomorphism \(\theta: \mathrm{U}' \to \mathrm{V}\) . Restricting \(\mathrm{U}'\) and V if necessary, we can assume that they are stable under H and that there exists an isomorphism \(u: \mathrm{T} \to \mathrm{V}\) compatible with the operations of H (Lemma 4). It now suffices to take \(\varphi = \theta^{-1} \circ u\) .

Recall (Differentiable and Analytic Manifolds, Results, 6.5.1) that if G is a Lie group, H a Lie subgroup of G and Y a manifold on which H operates on the left, we denote by \(G \times^{H}Y\) the quotient of the product manifold \(G \times Y\) by the right operation \(((g,y),h) \mapsto (gh,h^{-1}y)\) of H; this is a manifold on which the Lie group G operates naturally on the left; the projection \(G \times^{H}Y \to G/H\) is a bundle with fibre Y. Further, if Y is a finite dimensional vector space on which H operates linearly, \(G \times^{H}Y\) has a natural structure of vector G-bundle with base G/H (Differentiable and Analytic Manifolds, Results, 7.10.2).

Let G be a Lie group operating properly on the manifold X (General Topology, Chap. III, §4, no. 1, Def. 1) such that the law of operation \((g, x) \mapsto gx\) is of class \(C^{r}\) . Then, for every point x of X, the orbit Gx of x is a closed submanifold of X, isomorphic to the Lie homogeneous space \(G/G_{x}\) , where \(G_{x}\) is the fixer of x in G (cf.~Chap. III, §1, no. 7, Prop. 14 (ii), and General Topology, Chap. III, §4, no. 2, Prop. 4); this is a compact Lie group (loc. cit.).

PROPOSITION 6. Assume that the manifold X is paracompact; let x be a point of X, \(G_{x}\) its fixer. There exists a finite dimensional analytic linear representation \(\tau: G_{x} \to \mathbf{GL}(W)\) , and an open embedding \(\alpha: G \times ^{G_{x}} W \to X\) of class \(C^{r}\) , compatible with the operations of G, that maps the class of \((e,0) \in G \times W\) to x.

Put \(T = T_{x}(X)\) . Let W be a complementary subspace of \(T_{x}(Gx)\) in T, stable under \(G_{x}\) (for example, the orthogonal complement of \(T_{x}(Gx)\) with respect to a scalar product on T invariant under \(G_{x}\) ). On the other hand, let \(\varphi : T \to X\) be a morphism with the properties stated in Prop. 5 (relative to \(H = G_{x}\) ). Consider the morphism \(\lambda : G \times W \to X\) defined by \(\lambda(g, w) = g\varphi(w)\) . It induces by passage to the quotient a morphism \(\mu : G \times^{G_{x}}W \to X\) of class \(C^{r}\) , compatible with the operations of G, that maps the class z of \((e, 0)\) to x.

We show that \(\mu\) is étale at the point \(z\) . We have

\[
\begin{array}{c} \dim (\mathrm{G} \times {} ^ {\mathrm{G} _ {x}} \mathrm{W}) = \dim (\mathrm{G}) + \dim (\mathrm{W}) - \dim (\mathrm{G} _ {x}) \\ = \dim (\mathrm{G} x) + \dim (\mathrm{W}) = \dim (\mathrm{T}), \end{array}
\]

so it suffices to show that \(\mu\) is submersive at \(z\) , or equivalently that \(\lambda\) is submersive at \((e,0)\) . But, the tangent map \(\mathrm{T}_{(e,0)}(\lambda):\mathrm{T}_e(\mathrm{G})\oplus \mathrm{W}\to \mathrm{T}\) is equal to \(\mathrm{T}_e(\rho (x)) + i\) , where \(\rho (x)\) is the orbital map \(g\mapsto gx\) and \(i\) the canonical injection from W to T; since \(\operatorname {Im}\operatorname {T}_e(\rho (x)) = \operatorname {T}_x(\operatorname {G}x)\) , the map \(\mathrm{T}_{(e,0)}(\lambda)\) is surjective, and \(\mu\) is étale at \(z\) .

We are going to show that there exists an open neighbourhood \(\Omega\) of \(Gz\) in \(G \times ^{G_x}W\) , stable under \(G\) , such that \(\mu\) induces an isomorphism from \(\Omega\) onto an open subset of \(X\) . This will imply the proposition: indeed, the inverse image of \(\Omega\) in \(G \times W\) is stable under \(G\) , and hence is of the form \(G \times B\) , where \(B\) is an open subset of \(W\) containing the origin and stable under \(G_x\) ; restricting \(\Omega\) if necessary, we can assume that there exists an isomorphism \(u: W \to B\) , compatible with the operations of \(G_x\) (Lemma 4). It is clear that the composite morphism \(\alpha: G \times ^{G_x}W \xrightarrow{(\text{Id}, u)} G \times ^{G_x}B \xrightarrow{\mu} X\) satisfies the conditions in the statement of the proposition.

Thus, the proposition is a consequence of the following lemma:

Lemma 5. Let Z be a separated manifold of class \(C^{r}\) , equipped with a law of left operation \(m: G \times Z \to Z\) of class \(C^{r}\) , and \(\mu: Z \to X\) a morphism (of class \(C^{r}\) ) compatible with the operations of G. Let z be a point of Z, and \(x = \mu(z)\) . Assume that \(\mu\) is étale at z, and that the fixer of z in G is equal to the fixer \(G_{x}\) of x. Then, there exists an open neighbourhood \(\Omega\) of the orbit Gz, stable under G, such that \(\mu\) induces an isomorphism from \(\Omega\) onto an open subset of X.

Since \(\mu\) is compatible with the operations of G, it is étale at every point of Gz; since the canonical map \(G/G_{x} \rightarrow Gx\) is a homeomorphism, so is the map from Gz to Gx induced by \(\mu\) . Hence, it follows from Prop. 2 of no. 1 that there exists an open neighbourhood U of Gz in Z such that \(\mu\) induces an open embedding of U into X.

Since G operates properly on X, there exists an open neighbourhood V of x and a compact subset K of G such that \(gV \cap V = \varnothing\) for \(g \notin K\) (General Topology, Chap. III, §4, no. 4, Prop. 7); in particular, \(e \in K\) . The set \(W_{1}\) of points \(y \in Z\) such that \(Ky \subset U\) is open in Z: indeed, \(Z - W_{1}\) is the image of the closed set \((\mathrm{K} \times \mathrm{Z}) - m^{-1}(\mathrm{U})\) under the proper projection \(pr_{2}: K \times Z \to Z\) . Put \(W = W_{1} \cap \mu^{-1}(V)\) ; this is an open subset of Z, containing z, and satisfying the following conditions:

\begin{enumerate}
\def\labelenumi{(\roman{enumi})}
\item
  KW \(\subset\) U, and in particular W \(\subset\) U;
\item
  \(\mu(\mathrm{W}) \subset \mathrm{V}.\)
\end{enumerate}

Put \(\Omega = \mathrm{GW}\) and consider the restriction of \(\mu\) to \(\Omega\) . This is an étale morphism, since every point of \(\Omega\) is conjugate under \(\mathbf{G}\) to a point of \(\mathbf{U}\) . We show that it is injective: let \(g, h\) in \(\mathbf{G}\) and \(u, v\) in \(\mathbf{W}\) be such that \(\mu(gu) = \mu(hv)\) . Put \(k = g^{-1}h\) ; then \(\mu(u) = k\mu(v)\) , so \(k \in \mathbf{K}\) by (ii). But \(kv\) and \(u\) belong to \(\mathbf{U}\) by (i); thus, \(u = kv\) because the restriction of \(\mu\) to \(\mathbf{V}\) is injective, so \(gu = hv\) . Hence, the restriction of \(\mu\) to \(\Omega\) is injective, and consequently (Differentiable and Analytic Manifolds, Results, 5.7.8) is an isomorphism onto an open submanifold of \(\mathbf{X}\) , which completes the proof.

Under the conditions of Prop. 6, the image of \(\alpha\) is an open neighbourhood T of the orbit A of \(x\) , equipped with the structure of vector bundle with base A, for which the zero section is the orbit A itself. Such a neighbourhood is called a linear tube (around the orbit in question). For each point \(a \in \mathrm{A}\) , the fibre \(\mathrm{Y}_a\) of this vector bundle is a submanifold of X, stable under the fixer \(\mathrm{G}_a\) of \(a\) , and such that the morphism from \(\mathrm{G} \times^{\mathrm{G}_a}\mathrm{Y}_a\) to X that maps the class of \((g,y) \in \mathrm{G} \times \mathrm{Y}_a\) to \(gy \in \mathrm{X}\) induces a morphism of class \(\mathrm{C}^r\) from \(\mathrm{G} \times^{\mathrm{G}_a}\mathrm{Y}_a\) to T. Then \(\mathrm{Y}_a\) is said to be the transversal at \(a\) of the tube T. We remark that the tangent space at \(a\) of \(\mathrm{Y}_a\) is canonically isomorphic to \(\mathrm{Y}_a\) and that it is a complement of \(\mathrm{T}_a(\mathrm{A})\) in \(\mathrm{T}_a(\mathrm{X})\) ; thus, the vector bundle T with base A is canonically isomorphic to the normal bundle of A in X (Differentiable and Analytic Manifolds, Results, 8.1.3).

